\documentclass[12pt,a4paper]{article}

% ----------------------------------------------------------------------
% Pacotes
% ----------------------------------------------------------------------
\usepackage[utf8]{inputenc}
\usepackage[T1]{fontenc}
\usepackage[brazil]{babel}
\usepackage[a4paper,margin=2.5cm]{geometry}
\usepackage{booktabs}
\usepackage{tabularx}
\usepackage{enumitem}
\usepackage{xcolor}
\usepackage{titlesec}
\usepackage[hidelinks]{hyperref}

% ----------------------------------------------------------------------
% Aparência
% ----------------------------------------------------------------------
\definecolor{azul}{RGB}{13,110,253}
\titleformat{\section}{\large\bfseries\color{azul}}{\thesection}{1em}{}
\titleformat{\subsection}{\normalsize\bfseries}{\thesubsection}{1em}{}
\setlength{\parindent}{0pt}
\setlength{\parskip}{6pt}
\newcommand{\code}[1]{\texttt{#1}}

% ----------------------------------------------------------------------
% Título
% ----------------------------------------------------------------------
\title{\textbf{Sprint 5 --- Relação entre Implementação,\\ Arquitetura e Modelagem}\\[4pt]
\large Salvou CRM}
\author{Engenharia de Software II}
\date{\today}

\begin{document}
\maketitle

\begin{flushleft}
\textbf{Projeto:} Salvou CRM \\
\textbf{Repositório:} \url{https://github.com/llagogabriel/salvou-crm} \\
\textbf{Branch:} \texttt{Verificações}
\end{flushleft}

\hrulefill

% ----------------------------------------------------------------------
\section{Arquitetura adotada}

O Salvou CRM segue uma \textbf{arquitetura em camadas} baseada no padrão
\textbf{MVC (Model--View--Controller)}, implementada com o \textbf{Spring Boot}.
A separação de responsabilidades em camadas é o que mantém o código organizado e
coerente com a arquitetura planejada nas sprints anteriores.

\begin{center}
\renewcommand{\arraystretch}{1.3}
\begin{tabularx}{\textwidth}{l X}
\toprule
\textbf{Camada} & \textbf{Responsabilidade / Onde está no código} \\
\midrule
Apresentação (View) & Telas que o usuário vê e onde digita os dados --- templates
Thymeleaf em \code{resources/templates/} (\code{home.html}, \code{lista-clientes.html},
\code{cadastro-cliente.html}, \code{perfil-cliente.html}, \code{venda.html},
\code{agenda.html}) + Bootstrap/CSS. \\
Controle (Controller) & Recebe as requisições HTTP, orquestra a lógica e devolve a
tela --- \code{HomeController}, \code{ClienteController}, \code{VendaController},
\code{AgendaController}. \\
Modelo / Domínio (Model) & Representa as entidades do negócio --- \code{Cliente},
\code{Venda}, \code{Servico}, \code{Agenda}. \\
Persistência (Repository) & Acesso e gravação no banco --- \code{ClienteRepository},
\code{VendaRepository}, \code{ServicoRepository}, \code{AgendaRepository}. \\
Banco de dados & Armazenamento físico --- banco \textbf{H2} em arquivo
(\code{data/salvouDB}). \\
\bottomrule
\end{tabularx}
\end{center}

\subsection{Princípios da arquitetura aplicados no código}
\begin{itemize}[leftmargin=*]
  \item \textbf{Injeção de dependência / Inversão de controle:} os controllers não
        criam os repositórios manualmente; o Spring injeta as dependências via
        \code{@Autowired}, reduzindo o acoplamento entre as camadas.
  \item \textbf{Camada de persistência abstraída:} os repositórios são interfaces que
        estendem \code{JpaRepository}. Não há SQL manual --- o Spring Data gera o CRUD
        e as consultas a partir do nome dos métodos (ex.: \code{existsByCpf},
        \code{findByConcluidoFalse}, \code{findByDataAgendaBetween}).
  \item \textbf{Separação View/Controller:} os controllers não montam HTML; preparam
        os dados (via \code{Model}) e delegam a renderização ao Thymeleaf. Assim,
        mudar a tela não afeta a lógica e vice-versa.
\end{itemize}

% ----------------------------------------------------------------------
\section{Modelagem de classes e como foi implementada}

As classes de domínio definidas na modelagem (diagrama de classes) foram
implementadas como \textbf{entidades JPA} (anotação \code{@Entity}), em que cada
classe vira uma tabela e cada atributo vira uma coluna. Os relacionamentos do
diagrama foram traduzidos em anotações de mapeamento.

\subsection{Entidades}
\begin{itemize}[leftmargin=*]
  \item \textbf{Cliente} --- \code{id}, \code{nome}, \code{telefone}, \code{cpf}
        (único), \code{endereco}, \code{email}, \code{dataNascimento} e a lista de
        \code{vendas}.
  \item \textbf{Venda} --- \code{id}, \code{cliente}, lista de \code{servicos},
        \code{valorTotal}, \code{dataVenda}.
  \item \textbf{Servico} --- \code{id}, \code{nome}, \code{preco}.
  \item \textbf{Agenda} --- \code{id}, \code{descricao}, \code{tipo},
        \code{dataAgenda}, \code{concluido}, \code{observacoes} e um \code{cliente}
        (opcional).
\end{itemize}

\subsection{Relacionamentos (do diagrama para o código)}
\begin{center}
\renewcommand{\arraystretch}{1.3}
\begin{tabularx}{\textwidth}{X l X}
\toprule
\textbf{Relacionamento} & \textbf{Cardinalidade} & \textbf{Implementação} \\
\midrule
Cliente possui Vendas & 1 --- N & \code{@OneToMany} em \code{Cliente.vendas} /
\code{@ManyToOne} em \code{Venda.cliente} \\
Venda contém Serviços & N --- N & \code{@ManyToMany} em \code{Venda.servicos} com
tabela de junção \code{venda\_servicos} \\
Agenda referencia Cliente & N --- 1 (opcional) & \code{@ManyToOne} em
\code{Agenda.cliente} \\
\bottomrule
\end{tabularx}
\end{center}

Esses relacionamentos são exatamente os que o sistema usa em tempo de execução: ao
registrar uma venda, o vínculo \code{Venda~$\rightarrow$~Cliente} e
\code{Venda~$\rightarrow$~Serviços} é criado e gravado; ao abrir o perfil, a relação
\code{Cliente~$\rightarrow$~Vendas} é percorrida para somar o total gasto
(\code{getValorTotalVendas()}).

% ----------------------------------------------------------------------
\section{Rastreabilidade: Requisito $\rightarrow$ Classe $\rightarrow$ Arquitetura $\rightarrow$ Implementação}

A tabela mostra que a implementação mantém a relação entre o que foi pedido
(histórias de usuário), a modelagem (classes) e a arquitetura (camadas):

\begin{center}
\renewcommand{\arraystretch}{1.3}
\begin{tabularx}{\textwidth}{X l l}
\toprule
\textbf{Requisito / História} & \textbf{Modelo} & \textbf{Controle} \\
\midrule
Cadastrar/listar clientes; impedir CPF repetido & \code{Cliente} & \code{ClienteController} \\
Registrar venda e calcular o total & \code{Venda}, \code{Servico} & \code{VendaController} \\
Ver o total gasto por um cliente & \code{Cliente}, \code{Venda} & \code{ClienteController} \\
Organizar compromissos/agenda & \code{Agenda} & \code{AgendaController} \\
\bottomrule
\end{tabularx}
\end{center}

% ----------------------------------------------------------------------
\section{Coerência entre arquitetura, modelagem e implementação}

\begin{itemize}[leftmargin=*]
  \item \textbf{A arquitetura foi respeitada:} cada requisição segue o caminho
        previsto \emph{View $\rightarrow$ Controller $\rightarrow$ Repository
        $\rightarrow$ Banco $\rightarrow$ View}, sem misturar responsabilidades
        (a tela não acessa o banco diretamente; quem faz isso é o repositório,
        chamado pelo controller).
  \item \textbf{A modelagem foi seguida fielmente:} classes, atributos e
        relacionamentos do diagrama correspondem 1:1 às entidades JPA implementadas.
  \item \textbf{O escopo do MVP foi atendido:} as funcionalidades centrais (clientes,
        vendas e agenda) estão implementadas e funcionando de ponta a ponta, com
        persistência real.
\end{itemize}

Em resumo, a implementação não é um código solto: ela é a \textbf{materialização
direta} da modelagem de classes dentro da arquitetura MVC em camadas que o grupo
definiu nas sprints anteriores.

\end{document}
